% =====================================================================
%  Makra pro značení (převzato z 4EK216 makra.tex, převedeno do UTF-8, bez mtpro2)
%  \mbf → \mathbf, \what → \widehat, \wtilde → \widetilde, \upbeta → \beta
%  Tučná řečtina je stojatá přes \symbfup (unicode-math, xelatex); balíček bm se nepoužívá.
% =====================================================================

% Tučná písmena (vektory, matice)
\renewcommand{\aa}{\ensuremath{\mathbf{a}}}%
\newcommand{\bb}{\ensuremath{\mathbf{b}}}%
\newcommand{\cc}{\ensuremath{\mathbf{c}}}%
\newcommand{\dd}{\ensuremath{\mathbf{d}}}%
\newcommand{\ee}{\ensuremath{\mathbf{e}}}%
\newcommand{\ii}{\ensuremath{\mathbf{i}}}%
\newcommand{\pp}{\ensuremath{\mathbf{p}}}%
\newcommand{\qq}{\ensuremath{\mathbf{q}}}%
\newcommand{\rr}{\ensuremath{\mathbf{r}}}%
\renewcommand{\ss}{\ensuremath{\mathbf{s}}}%
\newcommand{\uu}{\ensuremath{\mathbf{u}}}%
\newcommand{\vv}{\ensuremath{\mathbf{v}}}%
\newcommand{\ww}{\ensuremath{\mathbf{w}}}%
\newcommand{\xx}{\ensuremath{\mathbf{x}}}%
\newcommand{\yy}{\ensuremath{\mathbf{y}}}%
\newcommand{\zz}{\ensuremath{\mathbf{z}}}%
\renewcommand{\AA}{\ensuremath{\mathbf{A}}}%
\newcommand{\BB}{\ensuremath{\mathbf{B}}}%
\newcommand{\CC}{\ensuremath{\mathbf{C}}}%
\newcommand{\DD}{\ensuremath{\mathbf{D}}}%
\newcommand{\EE}{\ensuremath{\mathbf{E}}}%
\newcommand{\II}{\ensuremath{\mathbf{I}}}%
\newcommand{\LL}{\ensuremath{\mathbf{L}}}%
\newcommand{\MM}{\ensuremath{\mathbf{M}}}%
\newcommand{\NN}{\ensuremath{\mathbf{N}}}%
\newcommand{\PP}{\ensuremath{\mathbf{P}}}%
\newcommand{\QQ}{\ensuremath{\mathbf{Q}}}%
\newcommand{\RR}{\ensuremath{\mathbf{R}}}%
\renewcommand{\SS}{\ensuremath{\mathbf{S}}}%
\newcommand{\TT}{\ensuremath{\mathbf{T}}}%
\newcommand{\UU}{\ensuremath{\mathbf{U}}}%
\newcommand{\VV}{\ensuremath{\mathbf{V}}}%
\newcommand{\XX}{\ensuremath{\mathbf{X}}}%
\newcommand{\YY}{\ensuremath{\mathbf{Y}}}%
\newcommand{\ZZ}{\ensuremath{\mathbf{Z}}}%
\newcommand{\nul}{\ensuremath{\mathbf{0}}}%  nulový vektor
\newcommand{\jedn}{\ensuremath{\mathbf{1}}}% vektor jedniček

% Řecká písmena a odhady
\newcommand{\eps}{\varepsilon}%
\newcommand{\aalpha}{\symbfup{\alpha}}%
\newcommand{\bbeta}{\symbfup{\beta}}%
\newcommand{\ddelta}{\symbfup{\delta}}%
\newcommand{\ggamma}{\symbfup{\gamma}}%
\newcommand{\ssigma}{\symbfup{\sigma}}%
\newcommand{\ttheta}{\symbfup{\theta}}%
\newcommand{\mmu}{\symbfup{\mu}}%
\newcommand{\ppi}{\symbfup{\pi}}%
\newcommand{\SSigma}{\symbfup{\Sigma}}%
\newcommand{\OOmega}{\symbfup{\Omega}}%
\newcommand{\eeps}{\symbfup{\varepsilon}}%
\newcommand{\betah}{\widehat{\beta}}%
\newcommand{\betat}{\widetilde{\beta}}%
\newcommand{\bbetah}{\widehat{\bbeta}}%
\newcommand{\bbetat}{\widetilde{\bbeta}}%
\newcommand{\ggammah}{\widehat{\ggamma}}%
\newcommand{\tthetah}{\widehat{\ttheta}}%
\newcommand{\tthetat}{\widetilde{\ttheta}}%
\newcommand{\SSigmah}{\widehat{\SSigma}}%
\newcommand{\OOmegah}{\widehat{\OOmega}}%
\newcommand{\Vhat}{\widehat{\mathbf{V}}}%
\newcommand{\uhat}{\widehat{u}}%
\newcommand{\uuhat}{\widehat{\mathbf{u}}}%
\newcommand{\yhat}{\widehat{y}}%
\newcommand{\yyhat}{\widehat{\mathbf{y}}}%

% Transpozice: stojaté bezpatkové T v horním indexu
\newcommand{\trans}{^{\mkern-1.5mu\mathsf{T}}}%
\newcommand{\tran}{^{\mkern-1.5mu\mathsf{T}\mkern-1.5mu}}%
\newcommand{\T}{^{\mkern-1.5mu\mathsf{T}\mkern-1.5mu}}%
\newcommand{\tra}{^{\mkern-1.5mu\mathsf{T}\mkern-1.5mu}}%
\newcommand{\inv}{^{-1}}%
\newcommand{\XTX}{\XX\trans\XX}%
\newcommand{\xb}{\xx\bbeta}%

% Operátory
% := piš jako \coloneq (centercolon z mathtools nefunguje s unicode-math)
\DeclarePairedDelimiter{\ceil}{\lceil}{\rceil}%
\DeclarePairedDelimiter{\floor}{\lfloor}{\rfloor}%
\DeclarePairedDelimiterX{\norm}[1]{\lVert}{\rVert}{#1}%
\let\Pr\relax
\DeclareMathOperator{\Pr}{Pr}%
\DeclareMathOperator{\prob}{Pr}%
\DeclareMathOperator{\ev}{E}%
\DeclareMathOperator{\var}{var}%
\DeclareMathOperator{\Avar}{Avar}%
\DeclareMathOperator{\sd}{sd}%
\DeclareMathOperator{\corr}{corr}%
\DeclareMathOperator{\cov}{cov}%
\DeclareMathOperator{\se}{se}% se(β̂ⱼ) stojatě: funkce odhadu
% samostatné se (hotová veličina, „se klasické“, „robustní se“) kurzívou jako t, p, F;
% platí se = se(β̂ⱼ) (konvence od 10. 11. 2026)
\newcommand{\sei}{\ensuremath{\mathit{se}}}%
\DeclareMathOperator{\diag}{diag}%
\DeclareMathOperator{\ind}{\perp\!\!\!\perp}%
\DeclareMathOperator{\logit}{logit}%
\DeclareMathOperator*{\argmax}{arg\,max}%
\DeclareMathOperator*{\argmin}{arg\,min}%
\DeclareMathOperator{\rank}{rank}%
\DeclareMathOperator{\tr}{tr}%
\DeclareMathOperator{\Span}{span}%
\DeclareMathOperator*{\plim}{plim}%
\newcommand{\Norm}{\mathcal{N}}% normální rozdělení (kaligrafické N), např. \Norm(0, 1)
\newcommand{\SST}{\mathrm{SST}}%
\newcommand{\SSE}{\mathrm{SSE}}%
\newcommand{\SSR}{\mathrm{SSR}}%
\newcommand{\abs}[1]{\ensuremath{|#1|}}%

% Šipky a konvergence
\newcommand{\sfi}{\ensuremath{\Rightarrow}}%
\newcommand{\sif}{\ensuremath{\Leftarrow}}%
\newcommand{\siff}{\ensuremath{\Leftrightarrow}}%
\newcommand{\lfi}{\ensuremath{\Longrightarrow}}%
\newcommand{\lif}{\ensuremath{\Longleftarrow}}%
\newcommand{\liff}{\ensuremath{\Longleftrightarrow}}%
\newcommand{\asim}{\ensuremath{\stackrel{\text{a}}{\sim}}}%
\newcommand{\pto}{\ensuremath{\stackrel{\text{p}}{\rightarrow}}}%
\newcommand{\dto}{\ensuremath{\stackrel{\text{d}}{\rightarrow}}}%

% Různé
\newcommand{\Rn}[1]{\ensuremath{\mathbb{R}^{#1}}}% reálné n-tice (koeficienty, čísla)
\newcommand{\En}[1]{\ensuremath{\mathbb{E}^{#1}}}% euklidovský prostor (geometrie: délka, kolmost, úhel)
\newcommand{\dif}{\,\mathrm{d}}%
\newcommand{\HO}{\ensuremath{\mathsf{H}_0}}%
\newcommand{\HI}{\ensuremath{\mathsf{H}_1}}%
\providecommand{\bm}[1]{\symbfup{#1}}% pro převzaté zdroje (\bm{\beta} → stojaté)
\newcommand{\aj}[1]{\textit{#1}}% anglický termín
\newcommand{\eng}[1]{(\textit{#1})}% anglický termín v závorce za českým
\newcommand{\pojem}[1]{\textbf{\textcolor{fisZelenaTmava}{#1}}}% nově zaváděný pojem
\newcommand{\kod}[1]{\texttt{#1}}%

% Odkaz na interaktivní materiál v Moodlu (jen slovně, bez URL)
\newcommand{\interaktivni}[1]{\textcolor{fisZelena}{\small$\blacktriangleright$}~Interaktivní materiál \textit{#1}}%
